% Figure 1 --- the same factor graph in its two conditioning patterns.
%
% This is the analogue of Wu & Tu's Figure 1 and it appears at the same point in the paper.
% The claim it carries is that encoding and decoding are two conditioning patterns of ONE
% model, so the two panels must be drawn with the same nodes, the same factor squares and the
% same parameters -- only the shading (what is observed), the arc directions and the frozen
% box may differ. Do not "improve" one panel without the other.
%
% Conventions: circles are variables, filled squares are factors, shaded circles are observed.
\begin{tikzpicture}[
  every node/.style={inner sep=1pt},
  var/.style={circle,draw,minimum size=13pt,font=\scriptsize},
  obs/.style={var,fill=black!14},
  fac/.style={rectangle,draw,fill=black!70,minimum size=4.2pt,inner sep=0pt},
  edge/.style={draw,black!55,line width=0.35pt},
  x=13.5mm, y=8.2mm,
]

% ---------------------------------------------------------------- panel (a): encoder
\begin{scope}
  \foreach \i in {1,2,3} {
    \node[obs]  (wa\i) at (\i,0)   {$w_\i$};
    \node[var]  (za\i) at (\i,1.5) {$Z_\i$};
    \node[fac]  (sa\i) at (\i,0.75) {};
    \draw[edge] (wa\i) -- (sa\i) -- (za\i);
  }
  % ternary factors, undirected: one per unordered pair, drawn above the labels
  \node[var,minimum size=11pt] (ha12) at (1.5,3.0) {$H$};
  \node[var,minimum size=11pt] (ha23) at (2.5,3.0) {$H$};
  \node[fac] (ta12) at (1.5,2.25) {};
  \node[fac] (ta23) at (2.5,2.25) {};
  \draw[edge] (za1) -- (ta12) -- (za2);
  \draw[edge] (ta12) -- (ha12);
  \draw[edge] (za2) -- (ta23) -- (za3);
  \draw[edge] (ta23) -- (ha23);
  \node[font=\scriptsize] at (2,-0.75) {(a) encoder};
\end{scope}

% ---------------------------------------------------------------- panel (b): decoder
\begin{scope}[xshift=48mm]
  % frozen prefix box
  \node[draw,dashed,black!55,rounded corners=1.5pt,fit={(0.60,-0.42) (2.40,2.02)},
        inner sep=0pt] (frozen) {};
  \node[font=\scriptsize,black!55,anchor=north] at (1.5,-0.44) {frozen prefix};

  \foreach \i in {1,2} { \node[obs] (wb\i) at (\i,0) {$w_\i$}; }
  \node[var] (wb3) at (3,0) {$W_3$};
  \foreach \i in {1,2,3} {
    \node[var] (zb\i) at (\i,1.5) {$Z_\i$};
    \node[fac] (sb\i) at (\i,0.75) {};
    \draw[edge] (wb\i) -- (sb\i) -- (zb\i);
  }
  % arcs reach backwards only: both heads of slot 3 sit on slot 3
  \node[var,minimum size=11pt] (hb1) at (2.30,3.05) {$H_3^{(c)}$};
  \node[fac] (tb1) at (2.0,2.25) {};
  \node[fac] (tb2) at (2.7,2.25) {};
  \draw[edge] (zb1) to[out=60,in=200] (tb1);
  \draw[edge] (tb1) -- (zb3);
  \draw[edge] (tb1) -- (hb1);
  \draw[edge] (zb2) to[out=70,in=210] (tb2);
  \draw[edge] (tb2) -- (zb3);
  \draw[edge] (tb2) -- (hb1);
  \node[font=\scriptsize] at (2,-1.15) {(b) decoder, $t=3$};
\end{scope}

\end{tikzpicture}
